\section{The Judgment at the Door}

The verse that carries the doctrine is one everybody knows and almost nobody reads to the end of the sentence.

``And as it is appointed unto men once to die, and after this the judgment: So also Christ was offered once to exhaust the sins of many. The second time he shall appear without sin to them that expect him unto salvation.''\endnote{Hebrews 9:27--28, Douay--Rheims, read 2026-07-26 in the tracked Challoner verse text. The Douay's ``to exhaust the sins of many'' renders \latin{ad multorum exhaurienda peccata}; the sense is that the sins are drawn off or taken away. The observation offered here concerns the syntax of the comparison (\latin{et quemadmodum\ldots{} sic et Christus}), which is on the surface of the Latin and the English alike.}

The famous half is a subordinate clause. Its job in the argument is to supply a familiar once-ness against which to set an unfamiliar one: men die once, then are judged; \emph{so also} Christ was offered once. The writer of Hebrews is not, at that moment, teaching about the individual's post-mortem judgment. He is teaching about the non-repeatability of the sacrifice, and he reaches for the one thing in human experience that everyone agrees happens exactly once.

That does not weaken the clause. It strengthens it in a particular way, because a premise assumed in passing is a premise the author expects no argument about. When a writer says \emph{as everyone knows, X; so also Y}, the interesting evidence is in the \emph{as everyone knows}. What the verse discloses is not a doctrine being proposed but a doctrine already in possession: among the people this letter addressed, that a man dies once and is then judged was not news.

It is worth being explicit about what the clause does and does not fix. It fixes the order---judgment after death, not before. It fixes the singularity---one death per man, which is why the Catechism cites it against reincarnation.\endnote{CCC 1013: ``When `the single course of our earthly life' is completed, we shall not return to other earthly lives: `It is appointed for men to die once.' There is no `reincarnation' after death.'' Holy See English web text, read 2026-07-26; the internal quotation is from \emph{Lumen gentium} 48. The Catechism's use of Hebrews 9:27 here is for the singularity, not for the judgment.} It does not fix the interval between them, and it does not say the judgment is individual rather than general. Everything further comes from elsewhere.

\subsection{What is defined, and what is only named}

Here the distinction between a doctrine and its vocabulary has to be made carefully, because the two are usually run together.

The Catechism's teaching is unambiguous: ``Each man receives his eternal retribution in his immortal soul at the very moment of his death, in a particular judgment that refers his life to Christ: either entrance into the blessedness of heaven---through a purification or immediately---or immediate and everlasting damnation.''\endnote{CCC 1022, Holy See English web text, read 2026-07-26. Its footnotes cite, for the purification, the Second Council of Lyons (1274), DS 857--858; the Council of Florence (1439), DS 1304--1306; and Trent (1563), DS 1820; for the immediate entrance, \emph{Benedictus Deus}, DS 1000--1001, and John XXII's \latin{Ne super his}, DS 990; and for the damnation, \emph{Benedictus Deus}, DS 1002. The appended sentence is from John of the Cross, \emph{Dichos} 64: ``At the evening of life, we shall be judged on our love.'' The essay read DS 990--991 and 1000--1002 directly; it did not read DS 857--858, 1304--1306, or 1820 at their own loci, and therefore relies on the Catechism's citation of them rather than asserting their contents.}

Now look at what the cited authorities actually say. \emph{Benedictus Deus} defines that the purified souls are in heaven and see the divine essence \latin{mox post mortem suam}---immediately after their death---and that the souls departing in actual mortal sin descend to hell \latin{mox post mortem suam}. That is a definition about \emph{timing and outcome}. In neither clause, nor in John XXII's declaration two years earlier, does the phrase \emph{particular judgment} occur.\endnote{The negative observation is bounded and should be read as such: within the texts examined for this essay---DS 990--991 and DS 1000--1002 in the Latin Enchiridion presentation read 2026-07-26---no equivalent of \latin{iudicium particulare} appears; DS 1002 uses \latin{in die iudicii} of the general judgment. The essay has not searched the wider magisterial corpus for the phrase and asserts nothing about where it first appears in an official text. The claim made in the body is only that the authorities the Catechism cites at CCC 1022 for the immediacy define an immediacy, and that the term is supplied by theology.} The substance is defined. The name is the theologians', and a very good name it is: if retribution is settled at death and the general judgment is still to come, then something judicial and individual has occurred, and it needs a word.

Why insist on the distinction? Because the name carries a picture, and the picture is not defined. A tribunal, an accounting, a reading-out of the record, an advocate and an accuser: all of that is imaginative furniture, some of it scriptural in origin and all of it inference. What is guaranteed is that at death the case is closed and the verdict is in force. How that presents itself to the one judged is not among the things the Church has told us, and the 1979 caution against ``arbitrary imaginative representations'' applies to this room as much as to any other.

There is also a small formal puzzle worth noticing, which turns out to be harmless. If the instant of death cannot be located---and the second section of this essay argued from Augustine that it cannot---then the instant of judgment cannot be located either, since the Catechism places it ``at the very moment of his death.'' But nothing follows from this except that the judgment is not an event within the world's time, measurable against other events in it. That was never a claim anyone made. The unlocatability of the moment is a fact about the observer's access to the boundary, not a defect in the boundary; the judgment is not late because we cannot time it.\endnote{This paragraph is project synthesis: an editorial reconciliation of CCC 1022's ``at the very moment of his death'' with the Augustinian argument reported in the second section, attributed to no source. It claims only that the two are compatible, not that any authority has related them.}

\subsection{The end of the time of merit}

The other half of the doctrine is the one that gives the whole series its urgency, and it is stated in the Catechism at its plainest: ``Death puts an end to human life as the time open to either accepting or rejecting the divine grace manifested in Christ.''\endnote{CCC 1021, Holy See English web text, read 2026-07-26; the paragraph's footnoted warrant is 2 Timothy 1:9--10. The same paragraph observes that ``the New Testament speaks of judgment primarily in its aspect of the final encounter with Christ in his second coming, but also repeatedly affirms that each will be rewarded immediately after death in accordance with his works and faith''---a candid acknowledgment that the scriptural centre of gravity is the general judgment and that the immediate one is attested more diffusely.} Compare the Catechism's earlier and more affecting version of the same claim: ``Death is the end of man's earthly pilgrimage, of the time of grace and mercy which God offers him so as to work out his earthly life in keeping with the divine plan, and to decide his ultimate destiny.''\endnote{CCC 1013, read as above.}

The scholastic version is blunter, and appears where one would least look for it: in the middle of Aquinas's argument that the immortality of the soul is no substitute for the resurrection. Explaining why Paul says that if Christ is not risen, then ``they also that are fallen asleep in Christ, are perished,'' he supplies the missing premise: they have perished ``because in the other life there are no meritorious works.''\endnote{Thomas Aquinas, \emph{Super I Cor.}, cap.~15, lect.~2: ``\latin{quia in alia vita nulla sunt opera meritoria}''; Latin read 2026-07-26 in the Corpus Thomisticum, in the same \latin{reportatio vulgata} noted above, and subject to the same textual caution. The English is a working gloss. The claim that the time of merit ends with death is common theology of long standing, not a definition; what the Church teaches definitively is the immediate retribution and the irrevocability of the outcome, from which the school claim is an inference about the mechanism.}

That the outcome is fixed is taught. \emph{Why} the will should be fixed at that moment---what it is about the separated condition, or about the completion of a life, that makes the orientation final---is not defined, and the schools have offered several accounts. This essay does not survey them, and a reader should treat any confident explanation of the mechanism, including the ones he has probably heard, as theology rather than as doctrine.\endnote{The essay deliberately does not enter the twentieth-century debates about whether the moment of death involves a final and fully lucid act of the person, nor the older scholastic accounts of the fixity of the separated will. Those discussions require sources this essay has not examined; the boundary is recorded here rather than crossed. What is asserted is only the distinction of weight: the fact of irrevocability is taught, the explanation is not defined.}

\subsection{What death does not sever}

One more negative is required, because the doctrine of an immediate individual judgment can be heard as making the dead unreachable, and the Church's practice says the opposite in the loudest possible way.

The Second Vatican Council states that the union between the living and ``the brethren who have gone to sleep in the peace of Christ is not in the least weakened or interrupted, but on the contrary, according to the perpetual faith of the Church, is strengthened by communication of spiritual goods,'' and that the Church ``from the very first ages of the Christian religion has cultivated with great piety the memory of the dead\ldots{} and offers suffrages for them.''\endnote{\emph{Lumen gentium} 49--50, Holy See English web text, read 2026-07-26; LG 50 cites 2 Maccabees 12:46. The theology of the purification presupposed by the practice of suffrage is treated in the second volume of this series and is not developed here.} And the 1979 letter, as noted, treats the funeral and the suffrage as themselves theological sources: any account of what happens after death which made the Church's prayers unintelligible would be refuted by the prayers.

So the judgment at death closes the individual's own account and does not close the communion. Those are different things, and the distinction is what allows a Catholic to hold simultaneously that his father's case is settled and that it is worth having a Mass said for him---a combination that looks incoherent to anyone who has collapsed judgment into finality of every kind.

That is as far as this volume goes into judgment. The event itself, its criterion, the two judgments and their relation, and the question of what could possibly be added at the end of the world to a verdict already in force, belong to \emph{Judgment}, the second volume of this series, which treats them at length.

What remains here is the part of the subject that the tradition thought most urgent and that modern Catholic writing has largely dropped: not what happens at death, but what a person is supposed to do about it beforehand.
